% !TeX root = ../../Book3.tex
\section{环域模量}
\label{b3:sec:0802}

模量量化两个边界分支之间的相对分离。它可以通过标准圆环定义，也可以不用求映射，直接由曲线族的长度与面积极值计算。后一种定义使比较和退化都可证明。

\subsection{环域模量}
\label{b3:subsec:080201}

\begin{definition}\label{b3:def:08-annulus-modulus}
若二连通区域 \(U\) 等价于 \(A(r,R)\)，定义\term[huanyu-moliang]{环域模量}[modulus of an annulus]为
\[
 \operatorname{mod}U=\frac1{2\pi}\log\frac Rr.
\]
对去心盘与去心平面记模量为 \(+\infty\)，并另记其退化类型。
\end{definition}
\symidx{\(\operatorname{mod}U\)：环域模量}
宽圆环的模量大，薄圆环的模量小；这与径向连接两边界的曲线“更长”的直觉一致。\(1/(2\pi)\) 来自绕行周长的归一化，不可在后面的矩形计算中随意删去。

\subsection{共形不变性}
\label{b3:subsec:080202}

\begin{proposition}\label{b3:prop:08-modulus-invariance}
两个有限圆环双全纯等价，当且仅当它们半径比相同。
\end{proposition}
\begin{proof}
必要性用本章第三节将完整证明的两项事实：极值长度在双全纯变换下不变，且连接圆环两端的曲线族的极值长度为 \((2\pi)^{-1}\log(R/r)\)。边界端由区域内部的趋近方式刻画，所以双全纯映射将此族对应到目标连接族，允许交换两端。充分性由缩放或倒数直接构造。
\end{proof}
这里证明只向前引用本章独立计算，不借模量不变性来证明极值长度不变性，因而没有循环。分类定理中的“唯一半径”也从此得到。

\subsection{单调性}
\label{b3:subsec:080203}

设 \(U\subset V\) 均为非退化二连通域，并且 \(U\) 中绕孔一周的曲线在 \(V\) 仍绕孔一周，称该包含为本质包含。由绕行曲线族的包含与下一节的倒数公式，有
\[
 \operatorname{mod}U\le\operatorname{mod}V.
\]
本质性必不可少。一个很小的环可以嵌在大环的局部圆盘内，它的孔在大环中可以被填掉；此时两者的模量没有由集合包含产生的固定关系。

\ProseParagraphBreak
具体地，\(A(r_2,R_2)\subset A(r_1,R_1)\) 若 \(r_1\le r_2<R_2\le R_1\)，中心孔相同，确为本质包含，模量不等式就是半径比的比较。对于非圆边界，只要能把区域夹在两个本质同心圆环之间，也可取得无需显式映射的上下界。

\begin{proposition}[串联环域的模量比较]\label{b3:prop:08-series-annuli}
设若干互不相交的二连通子域 \(U_1,\ldots,U_N\) 在环域 \(U\) 的两端之间依次分隔，且每条连接 \(U\) 两端的曲线都包含依次穿过各 \(U_j\) 的连接子曲线。则
\[
 \operatorname{mod}U\ge\sum_{j=1}^N\operatorname{mod}U_j.
\]
\end{proposition}
\begin{proof}
记 \(m_j=\operatorname{mod}U_j\)。由极值长度定义，在 \(U_j\) 上选密度 \(\rho_j\)，面积为 \(A_j>0\)，连接族的最短长度为 \(L_j>0\)，使 \(L_j^2/A_j\) 任意接近 \(m_j\)。把 \(\rho_j\) 乘 \(L_j/A_j\)，在其余区域置0，得到 \(U\) 上密度。每条总连接曲线长度至少为 \(\sum_j L_j^2/A_j\)，而总面积恰为这个和。因此总族极值长度至少为该和。逐一逼近各上确界便得所述式。零模量或极端情况先取有限近似，同样成立。
\end{proof}
同心圆环 \(A(r,s)\) 和 \(A(s,R)\) 在 \(A(r,R)\) 中串联，等号成为对数恒等式 \(\log(R/r)=\log(s/r)+\log(R/s)\)。中间圆周不属两个开环，但零面积的分隔曲线不改变所用长度积分。这项比较揭示细长瓶颈可以相加：只要每条总连接曲线必须穿过各段，局部阻碍便累积。若子环并列放置而曲线可以择一通过，串联假设失败，不能照搬求和。

\ProseParagraphBreak
可以给非本质包含一个数值反例。固定大环 \(V=A(1,3)\)，在点2附近放一个外半径为 \(\varepsilon<1/2\) 的小环 \(U=A_2(\varepsilon e^{-2\pi M},\varepsilon)\)，即把同心环平移到中心2。它完全包含于 \(V\)，模量恰为任意指定 \(M>0\)。当 \(M>\log3/(2\pi)\) 时，模量比大环还大；把半径比取接近1，又可任意小。小环的孔位于 \(V\) 内部，绕小孔的圈在 \(V\) 可收缩，所以该包含不是本质包含。由集合大小直接判断模量大小，会在此例失败。

\ProseParagraphBreak
串联比较也可使用极值密度的归一化解释。每段给一个最短连接长度1、面积接近 \(1/m_j\) 的密度。将其乘 \(m_j\)，长度贡献成为 \(m_j\)，面积贡献也成为 \(m_j\)。总路径必须付出各段长度之和，而总面积只按各段相加，因此总商至少为总模量。这相当于电阻串联，但证明并不依赖物理比喻；不相交支集与必须逐段穿过的条件已经给出全部数学依据。

\subsection{退化}
\label{b3:subsec:080204}

当 \(r\downarrow0\) 而 \(R\) 固定时，有限圆环模量趋于无穷，内端收缩为去心盘的点端。若 \(R/r\downarrow1\)，模量趋于0，径向厚度相对周长消失。两种极限分别对应长圆柱与薄圆柱；它们不应混称为同一种“退化”。

\ProseParagraphBreak
若同时令 \(r\downarrow0,R\uparrow\infty\)，局部区域穷竭为去心平面，模量也趋于无穷。与固定 \(R\) 的退化相比，第二端也走向无穷远，极限类型因此不同。只保存一个实数参数而不保存区域的观察点和端部数据，会在退化极限中丢失信息。
